% Recuperación aditiva del propietario N57, con enlace a las pruebas
% de cilindros y firma que ya están incluidas en el manuscrito.
\subsection{Cocycle de retenue et relèvement entier de la signature régionale}
\label{carry27:seccion}

La signature régionale conserve simultanément le résidu et le quotient.
La construction suivante rend explicite la loi algébrique de cette
dernière donnée. APP fournit les évaluations entières sur ses deux
feuillets ; TRIT conserve l'orientation et le régime ; le transport TPK
produit les bandes régionales et met à jour sa mémoire. Dans une
colonne de ces bandes, la division euclidienne détermine la retenue.
La construction utilise les représentants \(\{0,1,2\}\) de
\(\mathbb F_3\) ; la représentation orientée du TRIT conserve
séparément son signe et son quotient de transport.

Soient \((p,e,f)\in\{0,1,2\}^3\) les entrées d'une colonne régionale. Dans cette section, la lettre \(e\) désigne l'entrée ternaire de la région de propagation. Nous définissons
\begin{equation}
\begin{aligned}
q&=[p+e+f]_3,& c&=(p+e+f-q)/3,\\
a&=[p+e-f]_3,& h&=(p+e-f-a)/3,
\end{aligned}
\label{carry27:division}
\end{equation}
où \([z]_3\in\{0,1,2\}\) est le reste de tout entier \(z\).
La variable \(a\) est la coordonnée axiale de cette colonne. La
constante \(\alpha_{\rm HMT}\) conserve sa construction
dodécaphasique postérieure.

\begin{proposition}[Composition locale des retenues]
\label{carry27:cociclo}
Pour \(x,y,z\in\{0,1,2\}\), soient
\[
\kappa_3(x,y)=\frac{x+y-[x+y]_3}{3},\qquad
\beta_3(x,y)=\frac{x-y-[x-y]_3}{3}.
\]
Avec \(r=[p+e]_3\), les deux lois sont
\[
c=\kappa_3(p,e)+\kappa_3(r,f),\qquad
h=\kappa_3(p,e)+\beta_3(r,f).
\]
De plus,
\begin{equation}
\kappa_3(x,y)+\kappa_3([x+y]_3,z)
=\kappa_3(y,z)+\kappa_3(x,[y+z]_3).
\label{carry27:identidad}
\end{equation}
\end{proposition}

\begin{proof}
L'égalité \(p+e=r+3\kappa_3(p,e)\), ajoutée à
\(r+f=q+3\kappa_3(r,f)\), donne la première formule.
En soustrayant \(f\) et en utilisant
\(r-f=a+3\beta_3(r,f)\) on obtient la seconde.
Les deux membres de \eqref{carry27:identidad} sont
\((x+y+z-[x+y+z]_3)/3\). Par conséquent, la composition
préserve le même quotient quelle que soit l'association
de la somme.
\end{proof}

Si \(s_3:\mathbb F_3\to\mathbb Z\) est la section des
représentants, alors
\[
s_3(x)+s_3(y)-s_3(x+y)=3\kappa_3(x,y).
\]
Ceci est le cocycle de l'extension
\(0\to\mathbb Z\xrightarrow{\times3}\mathbb Z\to\mathbb F_3\to0\).
Son apparition provient du relèvement d'une opération résiduelle
à ses valeurs entières. La coordonnée visible et la mémoire du
multiple de trois restent liées par une identité exacte.

\Needspace{10\baselineskip} % S27_LAYOUT_TABLE
\begin{longtable}{ccc|rrrr}
\caption{Loi complète des \(27\) colonnes régionales.
Les valeurs de \(h\) sont présentées comme des entiers.}
\label{carry27:tabla}\\
\toprule
\(p\)&\(e\)&\(f\)&\(q\)&\(c\)&\(a\)&\(h\)\\
\midrule\endfirsthead
\toprule
\(p\)&\(e\)&\(f\)&\(q\)&\(c\)&\(a\)&\(h\)\\
\midrule\endhead
\bottomrule\endfoot
0&0&0&0&0&0&0\\
0&0&1&1&0&2&-1\\
0&0&2&2&0&1&-1\\
0&1&0&1&0&1&0\\
0&1&1&2&0&0&0\\
0&1&2&0&1&2&-1\\
0&2&0&2&0&2&0\\
0&2&1&0&1&1&0\\
0&2&2&1&1&0&0\\
1&0&0&1&0&1&0\\
1&0&1&2&0&0&0\\
1&0&2&0&1&2&-1\\
1&1&0&2&0&2&0\\
1&1&1&0&1&1&0\\
1&1&2&1&1&0&0\\
1&2&0&0&1&0&1\\
1&2&1&1&1&2&0\\
1&2&2&2&1&1&0\\
2&0&0&2&0&2&0\\
2&0&1&0&1&1&0\\
2&0&2&1&1&0&0\\
2&1&0&0&1&0&1\\
2&1&1&1&1&2&0\\
2&1&2&2&1&1&0\\
2&2&0&1&1&1&1\\
2&2&1&2&1&0&1\\
2&2&2&0&2&2&0\\
\end{longtable}

\begin{theorem}[Composante algébrique supplémentaire de la retenue]
\label{carry27:rango}
Sur \(\mathbb F_3\), les fonctions \(q,a\) sont linéaires en
\((p,e,f)\). Les réductions de \(c,h\) ont les représentations polynomiales suivantes :
\begin{align}
\bar c={}&2ef+2ef^2+2e^2f+2pf+2pf^2\notag\\
 &+2pe+pef+2pe^2+2p^2f+2p^2e,
\label{carry27:polinomio-c}\\
\bar h={}&2f^2+ef+2ef^2+e^2f+pf+2pf^2\notag\\
 &+2pe+2pef+2pe^2+p^2f+2p^2e.
\label{carry27:polinomio-h}
\end{align}
Si \(V\) est la matrice d'évaluation de \(1,p,e,f\) sur les
\(27\) colonnes, on a
\[
\operatorname{rank}_{\mathbb F_3}V=4,\qquad
\operatorname{rank}_{\mathbb F_3}(V\mid\bar c)=5,\qquad
\operatorname{rank}_{\mathbb F_3}(V\mid\bar h)=5.
\]
\end{theorem}

\begin{proof}
La substitution des \(27\) triplets du tableau dans les deux
polynômes reproduit \(\bar c,\bar h\). La représentation de
degré au plus deux en chaque variable est unique : un polynôme de
degré au plus deux en une variable qui s'annule aux trois
éléments de \(\mathbb F_3\) est nul ; l'application successive
de cet argument aux trois variables donne l'indépendance
des \(27\) monômes.

Les évaluations en \(0,(1,0,0),(0,1,0),(0,0,1)\) prouvent
l'indépendance de \(1,p,e,f\). La fonction \(\bar c\) vaut
zéro sur ces quatre entrées et vaut un sur \((1,1,1)\) ;
par conséquent, elle reste en dehors de son espace affine d'évaluation.
Pour \(\bar h\), les trois premières entrées imposent les
coefficients constant, de \(p\) et de \(e\) à zéro, et
\(\bar h(0,0,1)=2\) impose le coefficient de \(f\) à deux.
Cette fonction affine vaudrait un sur \((0,0,2)\), tandis
que \(\bar h(0,0,2)=2\). Chaque colonne supplémentaire augmente,
par conséquent, le rang d'une unité.
\end{proof}

La preuve précise quelle information apporte le quotient. La superposition résiduelle produit \(q,a\) ; l'évaluation entière incorpore en outre \(c,h\). Pour \(c\in\{0,1,2\}\), son résidu détermine la valeur entière, et pour \(h\in\{-1,0,1\}\), il la détermine conjointement avec ce choix déclaré de représentants. Les polynômes décrivent les réductions ; la mémoire intégrale conserve également ces relèvements.

\subsubsection{Transport des quotients régionaux}

Soit \(M\) une matrice entière de transport déjà fixée à une étape
du développement, et soient \(x_t^\chi\) ses états de ligne pour
\(\chi\in\{\pi,e,\varphi\}\). Pour chaque coordonnée, on effectue
la division
\[
y_t^\chi=x_t^\chi M,\qquad
u_t^\chi=[y_t^\chi]_3,\qquad
x_{t+1}^\chi=(y_t^\chi-u_t^\chi)/3.
\]
Définissons \(X_t^+=x_t^\pi+x_t^e+x_t^\varphi\) et
\(X_t^\sigma=x_t^\pi+x_t^e-x_t^\varphi\).
Les lois précédentes, appliquées colonne par colonne, donnent
\begin{equation}
\begin{aligned}
X_{t+1}^+&=\frac{X_t^+M-q_t}{3}-c_t,\\
X_{t+1}^\sigma&=\frac{X_t^\sigma M-a_t}{3}-h_t.
\end{aligned}
\label{carry27:transporte}
\end{equation}
En effet, la somme des trois résidus est \(q_t+3c_t\) ; la combinaison signée est \(a_t+3h_t\). Substituer les deux identités dans la division des états démontre \eqref{carry27:transporte}. L'énoncé a pour domaine les transports entiers déclarés. Les émissions régionales et leur continuité conservent les lecteurs coinductifs construits précédemment.

Comme évaluation complète sur une bande de six colonnes,
\[
u^\pi=021020,\quad u^e=001220,\quad u^\varphi=100210
\]
produisent \(q=122120\), \(c=000110\), \(a=222000\) et  
\(h=(-1,0,0,0,1,0)\). L'orientation  
\(\sigma=(1,1,-1)\) possède un résidu double  
\(2\sigma=(2,2,1)\) en \(\mathbb F_3^3\). Cette dernière  
égalité enregistre la convention de signe et se distingue des  
bandes de six colonnes.

\subsubsection{Mémoire de frontière et lecture de l'intervalle}

La loi locale de retenue s'accouple au raffinement cylindrique
déjà démontré. Un exemple montre la fonction de l'intervalle
complet :
\[
\left[\frac1{729},\frac2{729}\right)
\cap
\left[\frac2{1000},\frac3{1000}\right)
=
\left[\frac2{1000},\frac2{729}\right)\ne\varnothing.
\]
En revanche,  
\(\lfloor1000/729\rfloor=1\).  
La lecture de l'extrémité inférieure renvoie le bloc \(1\),  
alors que la compatibilité par intervalles admet également la  
cellule \(2\). L'état de frontière conserve les deux bornes  
nécessaires pour décider de cette compatibilité.

Pour une colonne, fixer \(q\) et \(a\) fixe
\[
f=2(q-a),\qquad p+e=q-f\quad\text{en }\mathbb F_3.
\]
Chacune des neuf fibres de \((q,a)\) contient exactement trois triplets : choisir \(p\) détermine \(e\) et laisse \(f\) fixe. Sur six colonnes, avant d'appliquer des restrictions supplémentaires, la liberté est le produit de ces six fibres et a pour cardinal \(3^6=729\). La fixation se rapporte aux alternatives d'une même frontière. En avançant en profondeur, le registre et sa mémoire changent selon le transport. Ainsi, la retenue locale, les bornes du cylindre et l'incidence régionale agissent conjointement sur la même histoire arithmétique.
